% !TeX root = part6.tex
\documentclass[../final.tex]{subfiles}

\begin{document}

% Продолжение по плану семинаров: тема 4, задача 2; сборник, задача 4.17.
\practice{6}{10.10.2026}{Средний квадрат смещения атома.
Флуктуации в низкоразмерных решётках}

\begingroup
\renewcommand{\thesection}{4.\arabic{section}}
\task[17]{Средний квадрат смещения атома в решётке}
\label{task:fksv_displacements}

\condition
Найти средний квадрат смещения атома в гармонической решётке
с помощью операторов рождения и уничтожения фононов и независимо
с помощью теоремы вириала. Исследовать сходимость результата при
$T=0$ и $T>0$ в размерностях $d=1,2,3$.

\solution

\subsection*{Нормальные моды и усреднение}

Продолжим квантование из предыдущего занятия. Сначала рассматриваем
одну скалярную ветвь одномерной цепочки из $N$ одинаковых атомов массы $m$.
Центр масс фиксирован, поэтому нулевая мода исключена.
Физическое волновое число обозначаем через $k$; в отличие от
безразмерного $k$ предыдущего занятия, здесь в экспоненте стоит $kna$.
Оператор смещения равен
\[
    \hat u_n=\frac1{\sqrt N}\sum_{k\ne0}
       \sqrt{\frac{\hbar}{2m\omega_k}}
       \left(\hat b_k e^{ikna}+\hat b_k^\dagger e^{-ikna}\right).
\]
В тепловом равновесии
\begin{align*}
    \langle\hat b_k^\dagger\hat b_{k'}\rangle
       &=\bar n_k\delta_{kk'},&
    \bar n_k&=\frac1{e^{\hbar\omega_k/(k_{\mathrm B}T)}-1},\\
    \langle\hat b_k\hat b_{k'}^\dagger\rangle
       &=(\bar n_k+1)\delta_{kk'},&
    \langle\hat b_k\hat b_{k'}\rangle
       &=\langle\hat b_k^\dagger\hat b_{k'}^\dagger\rangle=0.
\end{align*}
Среднее смещение равно нулю. При возведении $\hat u_n$ в квадрат
выживают два смешанных произведения, поэтому
\begin{align*}
    \langle u_n^2\rangle
       &=\frac1N\sum_{k\ne0}\frac{\hbar}{2m\omega_k}
           \bigl[(\bar n_k+1)+\bar n_k\bigr]\\
       &=\frac{\hbar}{2mN}\sum_{k\ne0}\frac1{\omega_k}
           \coth\frac{\hbar\omega_k}{2k_{\mathrm B}T}.
\end{align*}
Здесь использовано тождество $2\bar n_k+1=\coth(\hbar\omega_k/2k_{\mathrm B}T)$.
Результат не зависит от узла $n$ в силу трансляционной симметрии.
Слагаемое $1/2$ в энергии осциллятора отвечает за
\rconcept{нулевые колебания} и сохраняется даже при $T=0$.

\subsection*{Вывод через теорему вириала}

Выберем вещественные ортонормированные нормальные координаты $X_\nu$,
чтобы каждая независимая мода была обычным осциллятором:
\[
    H_\nu=\frac{P_\nu^2}{2m}+\frac{m\omega_\nu^2X_\nu^2}{2}.
\]
Для квадратичного потенциала теорема вириала даёт
\[
    \langle K_\nu\rangle=\langle V_\nu\rangle
       =\frac12\langle H_\nu\rangle,
    \qquad
    \langle H_\nu\rangle=\hbar\omega_\nu
       \left(\bar n_\nu+\frac12\right).
\]
Следовательно,
\[
    \frac{m\omega_\nu^2}{2}\langle X_\nu^2\rangle
       =\frac{\hbar\omega_\nu}{2}
          \left(\bar n_\nu+\frac12\right),
    \qquad
    \langle X_\nu^2\rangle
       =\frac{\hbar}{m\omega_\nu}
          \left(\bar n_\nu+\frac12\right).
\]
Ортогональность преобразования сохраняет сумму квадратов:
$\sum_n\langle u_n^2\rangle=\sum_\nu\langle X_\nu^2\rangle$.
Деля на $N$, вновь получаем найденный выше результат.
В комплексных координатах $Q_k$ тот же смысл имеет
$\langle Q_kQ_{-k}\rangle$; координаты $Q_k$ и $Q_{-k}$ нельзя
считать двумя независимыми вещественными осцилляторами.

\subsection*{Переход к произвольной размерности}

Для одноатомной решётки в $d$ измерениях векторное смещение
содержит сумму по поляризациям $s$. Нормируем векторы поляризации
условием $\boldsymbol e_{\boldsymbol k s}^*
\mathbin{\cdot}\boldsymbol e_{\boldsymbol k s}=1$. Тогда
\[
    \langle|\boldsymbol u|^2\rangle
       =\frac{\hbar}{2mN}\sum_{\boldsymbol k,s}'
          \frac1{\omega_{\boldsymbol k s}}
          \coth\frac{\hbar\omega_{\boldsymbol k s}}{2k_{\mathrm B}T}.
\]
Штрих исключает поступательные нулевые моды. Пусть $v_0=V_d/N$ ---
объём, площадь или длина на один атом. При больших размерах
\[
    \frac1N\sum_{\boldsymbol k}\longrightarrow
       \frac{v_0}{(2\pi)^d}\int d^dk.
\]
Дальнейший анализ проводим для
\rassumption{акустических ветвей с линейной дисперсией}
$\omega_{\boldsymbol k s}\simeq c_s k$.
Это существенно: для изгибной ветви свободной мембраны
с $\omega\propto k^2$ степени расходимости будут другими.

Обозначим $S_{d-1}=2\pi^{d/2}/\Gamma(d/2)$ --- площадь единичной
сферы в $d$ измерениях. Опасная длинноволновая часть интеграла равна
\[
    \langle|\boldsymbol u|^2\rangle_{\mathrm{IR}}
       =\frac{\hbar v_0 S_{d-1}}{2m(2\pi)^d}
          \sum_s\frac1{c_s}
          \int_{k_{\min}}^{k_*}
             k^{d-2}\coth\frac{\hbar c_sk}{2k_{\mathrm B}T}\,dk.
\]
Здесь $k_{\min}\sim2\pi/L$, а $k_*$ --- фиксированная граница
области линейной дисперсии. Конечная решётка сама задаёт
инфракрасное обрезание; расходимость возникает при $L\to\infty$.

\subsection*{Нулевая температура}

При $T=0$ гиперболический котангенс равен единице:
\[
    \langle|\boldsymbol u|^2\rangle_{T=0,\mathrm{IR}}
       \propto\int_{2\pi/L}^{k_*}k^{d-2}\,dk.
\]
Отсюда при $d=1$ имеем $\ln L$, а при $d=2,3$ интеграл
сходится на нижнем пределе. Для одной линейной ветви цепочки
с $v_0=a$ коэффициент логарифма равен
\[
    \langle u^2\rangle_{T=0}
       =\frac{\hbar a}{2\pi mc}\ln\frac{k_*L}{2\pi}+O(1).
\]
Верхняя граница зоны Бриллюэна конечна, поэтому ультрафиолетовая
часть решёточной модели дополнительной расходимости не создаёт.

\subsection*{Любая ненулевая температура}

При сколь угодно малой, но фиксированной $T>0$ найдутся моды
с $\hbar c_sk\ll k_{\mathrm B}T$. Для них
\[
    \coth\frac{\hbar c_sk}{2k_{\mathrm B}T}
       \simeq\frac{2k_{\mathrm B}T}{\hbar c_sk},
\]
и
\[
    \langle|\boldsymbol u|^2\rangle_{\mathrm{IR}}
       \simeq\frac{k_{\mathrm B}T v_0 S_{d-1}}{m(2\pi)^d}
          \sum_s\frac1{c_s^2}
          \int_{2\pi/L}^{k_*}k^{d-3}\,dk.
\]
Таким образом, в одном измерении дисперсия растёт как $TL$,
в двух --- как $T\ln L$, а в трёх остаётся конечной.
Для двумерной решётки коэффициент логарифма имеет вид
\[
    \langle|\boldsymbol u|^2\rangle_{\mathrm{IR}}
       =\frac{k_{\mathrm B}Tv_0}{2\pi m}
          \sum_s\frac1{c_s^2}\ln\frac{k_*L}{2\pi}+O(1).
\]
При низкой температуре верхний предел классической области следует
ограничить также масштабом $k_{\mathrm B}T/(\hbar c_s)$;
это меняет конечную часть, но не зависимость от $L$.

\subsection*{Точная проверка для классической конечной цепочки}

В высокотемпературном пределе исходная сумма даёт
\[
    \langle u^2\rangle
       =\frac{k_{\mathrm B}T}{mN}\sum_{j=1}^{N-1}\frac1{\omega_j^2},
    \qquad
    \omega_j^2=\frac{4\alpha}{m}\sin^2\frac{\pi j}{N}.
\]
Используем тригонометрическую сумму
\[
    \sum_{j=1}^{N-1}\frac1{\sin^2(\pi j/N)}=\frac{N^2-1}{3}.
\]
Её можно получить, дважды дифференцируя логарифм произведения
$\prod_{j=0}^{N-1}\sin(x+\pi j/N)=C_N\sin Nx$
и затем сравнивая постоянные члены при $x\to0$.
В результате
\[
    \langle u^2\rangle=\frac{k_{\mathrm B}T}{12\alpha}
       \left(N-\frac1N\right).
\]
Это явная проверка линейного роста с длиной цепочки $L=Na$.
Численный коэффициент определяется граничными условиями;
здесь цепочка периодическая и центр масс фиксирован.

\subsection*{Абсолютные и относительные смещения}

Расходимость абсолютного смещения не означает разрушения каждой
связи между соседями. Для расстояния $r$ в одномерной цепочке
\[
    \langle(u_{n+r/a}-u_n)^2\rangle
       =\frac{\hbar}{mN}\sum_{k\ne0}
          \frac{1-\cos kr}{\omega_k}
          \coth\frac{\hbar\omega_k}{2k_{\mathrm B}T}.
\]
При фиксированном $r$ множитель $1-\cos kr\simeq k^2r^2/2$
подавляет вклад самых длинных волн. В классической периодической
цепочке для $r=\ell a$ точный результат равен
\[
    \langle(u_{n+\ell}-u_n)^2\rangle
       =\frac{k_{\mathrm B}T}{\alpha}\ell
          \left(1-\frac\ell N\right),\qquad 0\le\ell\le N.
\]
Соседние атомы остаются относительно близки, но при увеличении
расстояния флуктуации накапливаются.
В двумерной гармонической модели при $T>0$ они растут
логарифмически с расстоянием: строгого дальнего трансляционного
порядка нет, хотя возможен квазидальний порядок.

\answer
\[
    \langle|\boldsymbol u|^2\rangle
       =\frac{\hbar}{2mN}\sum_{\boldsymbol k,s}'
          \frac1{\omega_{\boldsymbol k s}}
          \coth\frac{\hbar\omega_{\boldsymbol k s}}{2k_{\mathrm B}T}.
\]
Для линейных акустических ветвей зависимость от размера имеет вид
\[
\begin{array}{c|ccc}
    &d=1&d=2&d=3\\ \hline
    T=0&\ln L&\text{конечна}&\text{конечна}\\
    T>0&L&\ln L&\text{конечна}
\end{array}
\]
\rresult{Тепловые длинноволновые флуктуации расходятся в 1D и 2D.}
Вывод относится к термодинамическому пределу, фиксированному центру
масс и указанному закону дисперсии.

\endgroup
\end{document}
